\documentclass[aspectratio=169,10pt]{beamer}

\usepackage{kotex}
\usepackage{amsmath,amssymb}
\usepackage{booktabs}
\usepackage{tabularx}
\usepackage{array}
\usepackage{graphicx}
\usepackage{ragged2e}

\setmainhangulfont{AppleMyungjo}
\setsanshangulfont{Apple SD Gothic Neo}

\usetheme{Madrid}
\usecolortheme{default}
\setbeamertemplate{navigation symbols}{}
\setbeamertemplate{footline}{}
\setbeamertemplate{section in toc}[sections numbered]
\setbeamertemplate{subsection in toc}[subsections numbered]
\setbeamertemplate{itemize item}{\small$\blacksquare$}
\setbeamertemplate{itemize subitem}{\tiny$\blacktriangleright$}

\title{주행 상황 실시간 경고 시스템}
\subtitle{구간 정보 및 긴급차량 접근 기반 경고 컨셉\\전주기 시스템 설계·구현·검증 포트폴리오}
\author{Driving Alert Project}
\date{2026-03}

\begin{document}

\begin{frame}
  \titlepage
\end{frame}

\begin{frame}{목차}
\begin{columns}[T,onlytextwidth]
\column{0.52\textwidth}
\footnotesize
\tableofcontents

\column{0.46\textwidth}
\small
\begin{block}{발표 흐름}
\scriptsize
\begin{itemize}
  \item 전체 차량 기준선 $\rightarrow$ 다수 ECU 구조 $\rightarrow$ 네트워크 흐름
  \item \texttt{V2X} $\rightarrow$ \texttt{ADAS} $\rightarrow$ \texttt{CGW} $\rightarrow$ 진단/관측 계층
  \item 패널 $\rightarrow$ 요구사항 $\rightarrow$ 기능 정의 $\rightarrow$ 시험 전략
  \item \texttt{UT / IT / ST} $\rightarrow$ 실행 증빙 $\rightarrow$ 리포트 $\rightarrow$ 운영 자동화
  \item 정량 지표 구조 $\rightarrow$ appendix / 별첨
\end{itemize}
\end{block}
\end{columns}
\end{frame}

\AtBeginSubsection[]
{
  \begin{frame}{\insertsectionnumber.\insertsubsectionnumber\ \insertsubsectionhead}
  \footnotesize
  \tableofcontents[currentsection,currentsubsection,sectionstyle=show/shaded,subsectionstyle=show/show/hide]
  \end{frame}
}

\section{서론 및 문제 배경}
\subsection{문제 인식과 컨셉 도출}

\begin{frame}{프로젝트 출발점}
\begin{itemize}
  \item 출발점은 단순했다. 차량에는 이미 다양한 경고 기능이 존재하지만, 실제 주행에서는 여러 맥락이 동시에 들어오며 운전자는 우선순위가 정해진 최종 판단을 경험한다.
  \item 예를 들어 스쿨존 진입, 긴급차량 접근, 교차로 측방 위험이 겹치면 개별 기능은 각자 맞더라도 전체 경고는 쉽게 충돌한다.
  \item 그래서 본 프로젝트는 기존 차량 하드웨어 위에서 분산된 경고 기능을 소프트웨어 정책으로 다시 엮어, 새로운 주행 경고 컨셉으로 재구성하는 과제에서 시작되었다.
\end{itemize}
\end{frame}

\begin{frame}{문제 관찰}
\begin{itemize}
  \item 구간 경고와 긴급 경고가 동시에 유효해도 최종 경고가 일관되게 유지되지 않는 경우가 먼저 보였다.
  \item 경고 시작, 전환, 종료 과정에서는 표시 중복과 복귀 불안정이 반복되었고, 어떤 경고가 최종 상태여야 하는지가 명확하지 않았다.
  \item 결국 문제는 개별 기능 부족이 아니라, 전체 차량 관점에서 선택 규칙과 복귀 규칙이 분리돼 있었다는 점이었다.
\end{itemize}
\end{frame}

\begin{frame}{왜 통합 경고 컨셉이 필요했는가}
\begin{itemize}
  \item 스쿨존, 고속도로, 유도선은 각각 다른 구간 경고를 요구하지만, 운전자 입장에서는 모두 같은 주행 맥락 판단 문제다.
  \item 경찰차와 구급차 접근은 구간 경고보다 우선해야 하고, 동시 접근 시에도 하나의 선택 규칙이 필요하다.
  \item 교차로·합류구간 위험은 객체 위험과 긴급 이벤트를 함께 고려해야 하므로, 개별 기능 설명만으로는 시스템 전체를 설명할 수 없다.
\end{itemize}
\end{frame}

\begin{frame}{아이디어 도출}
\begin{block}{아이디어 요약}
\small
주행 중 운전자가 실제로 마주하는 장면은 기능별로 분리되어 나타나지 않는다. 구간 맥락은 \texttt{IVI/NAV}에서, 긴급 이벤트는 \texttt{V2X}에서, 충돌위험은 \texttt{ADAS}에서 들어오지만, 차량은 운전자에게 우선순위가 정해진 최종 경고 상태와 일관된 복귀 규칙을 제공해야 한다. 본 프로젝트는 이 분산된 입력을 같은 정책 위에서 해석하고, 기존 차량 구조 위에 소프트웨어 중심 경고 컨셉으로 재구성해 보자는 문제의식에서 출발하였다.
\end{block}
\end{frame}

\begin{frame}{컨셉 정의}
\begin{itemize}
  \item 본 컨셉은 구간 경고, 긴급차량 접근 경고, 교차로·합류구간 충돌위험 경고를 하나의 경고 정책으로 통합한다.
  \item 입력은 여러 개일 수 있지만, 최종 경고는 하나의 선택 규칙과 하나의 복귀 규칙을 통해 결정되도록 설계하였다.
  \item 결과적으로 운전자가 보는 출력은 앰비언트, 클러스터, 전면 표시, 오디오, 감속 보조 요청으로 분기되더라도 같은 경고 의미를 유지하도록 구성하였다.
\end{itemize}
\end{frame}

\subsection{연구 범위와 핵심 질문}
\begin{frame}{연구 범위}
\begin{itemize}
  \item 기본 차량 상태, 구간 정보, 긴급 이벤트, TTC 기반 객체 위험을 입력 범위로 정의하였다.
  \item 구현과 검증 범위는 CANoe SIL, CAN + Ethernet, CAPL 기반 논리 ECU 구조로 고정하였다.
  \item 실차 하드웨어, HIL, 상용 UDS 스택, 양산 timing 보증은 현재 사이클의 범위에서 제외하였다.
\end{itemize}
\end{frame}

\begin{frame}{핵심 연구 질문}
\begin{itemize}
  \item 어떤 입력이 들어왔는가보다, 최종적으로 어떤 경고가 선택되어야 하는가
  \item 경고가 잘 울리는가보다, 언제 안전하게 해제되고 복귀하는가
  \item 기능을 만들었는가보다, 설계·코드·시험·운영이 같은 기준선을 보는가
\end{itemize}
\end{frame}

\section{프로젝트 프로세스와 추적성 체계}
\subsection{문서 기반 개발 프로세스}

\begin{frame}{전체 수행 프로세스}
\begin{itemize}
  \item 프로젝트는 요구사항 정리에서 시작해 기능 정의, 네트워크/변수 설계, 구현, 시험, 운영 자동화로 이어지는 전체 프로세스로 진행되었다.
  \item 문서와 구현, 시험 자산을 따로 두지 않고 같은 기준선 안에서 갱신하는 방식이 핵심 운영 원칙이었다.
  \item 즉 본 프로젝트의 산출은 코드만이 아니라, 코드와 시험을 설명할 수 있는 문서 체계 자체였다.
\end{itemize}
\end{frame}

\begin{frame}{문서 체계와 단계별 산출물}
\begin{itemize}
  \item 00은 개요와 거버넌스 기준을 정리하는 상위 문서였다.
  \item 01은 무엇을 만들 것인가를, 03--0304는 어떻게 설계할 것인가를, 04는 구현 구조를 설명하였다.
  \item 05, 06, 07은 단위시험, 통합시험, 시스템시험을 통해 같은 기능을 서로 다른 수준에서 검증하도록 구성되었다.
\end{itemize}
\end{frame}

\begin{frame}{프로젝트 개요와 기준선 설정}
\begin{itemize}
  \item 프로젝트 개요 문서는 목적, 핵심 시나리오, 설계 규칙, 기대 가치를 한 페이지로 요약하는 출발점 역할을 하였다.
  \item 여기서 구간 경고, 긴급차량 접근, 교차로·합류 위험, 경고 충돌과 복귀를 핵심 시나리오로 고정하였다.
  \item 이후 세부 문서는 모두 이 개요 문서의 의미를 구체화하는 방식으로 연결되었다.
\end{itemize}
\end{frame}

\begin{frame}{요구사항 기준과 범위}
\begin{itemize}
  \item 요구사항은 구간 경고, 긴급 경고, 우선순위, 충돌위험, fail-safe, 강건성, 확장 ECU 연계까지 포함하도록 정리되었다.
  \item 특히 긴급 이벤트 종료, 1000ms 무갱신 해제, ETA/SourceID 동률 처리와 같은 규칙은 요구사항 수준에서 먼저 고정하였다.
  \item 이 단계에서 정한 규칙이 이후 기능 정의와 시험 기대값의 기준이 되었다.
\end{itemize}
\end{frame}

\begin{frame}{HARA와 안전 기준}
\begin{itemize}
  \item HARA는 구간 경고, 긴급 경고, fail-safe, 복귀 안정성 같은 위험 항목을 상위 안전 목표로 정리하는 역할을 하였다.
  \item 본 프로젝트는 ISO 26262 완전 준수를 선언하기보다, HARA를 통해 무엇이 안전 관점의 핵심인지 분명히 하는 데 초점을 두었다.
  \item 그 결과 위험 항목과 시험 시나리오를 같은 맥락에서 설명할 수 있는 구조가 만들어졌다.
\end{itemize}
\end{frame}

\subsection{시험 기준과 추적성}
\begin{frame}{시험 판정 기준과 Test Matrix}
\begin{itemize}
  \item Test Oracle은 경고가 언제 시작되고 무엇이 선택되며 언제 해제되어야 하는지를 판정하는 기준을 제공하였다.
  \item Test Matrix는 요구사항, HARA, 기능 정의, 시험 레벨을 하나의 표면에서 연결하는 상위 통제 문서로 사용되었다.
  \item 이를 통해 단위시험, 통합시험, 시스템시험이 서로 다른 시험이 아니라 하나의 추적 구조 위에 놓이게 되었다.
\end{itemize}
\end{frame}

\begin{frame}{추적성 체계}
\begin{block}{추적성 원칙}
\[
\text{Req} \rightarrow \text{Func} \rightarrow \text{Flow} \rightarrow \text{Comm} \rightarrow \text{Var} \rightarrow \text{Code} \rightarrow \text{Test}
\]
\end{block}
\begin{itemize}
  \item 각 문서는 위 화살표의 한 층을 맡고, 시험은 항상 상위 설계 의도까지 되돌아가서 설명될 수 있어야 했다.
  \item 이 구조는 ISO 26262 추적성과 V-cycle 검증 흐름을 설명하는 핵심 근거로 사용되었다.
\end{itemize}
\end{frame}

\section{시스템 아키텍처와 네트워크 설계}
\subsection{시스템 구조와 ECU 구성}

\begin{frame}{시스템 구조 개요}
\begin{itemize}
  \item 시스템은 입력 수집 계층, 판단 계층, 출력 계층, 검증 계층의 네 층으로 분리하였다.
  \item 입력 수집 계층은 차량 상태와 구간 맥락, 긴급 이벤트를 모으고, 판단 계층은 이를 단일 경고로 중재한다.
  \item 출력 계층은 경고를 운전자 채널로 반영하며, 검증 계층은 시나리오 입력과 관측을 분리해 유지한다.
\end{itemize}
\end{frame}

\begin{frame}{채택 아키텍처 구조}
\begin{itemize}
  \item 채택 구조는 \texttt{Domain CAN + ETH\_Backbone} 위에 중앙 경고 판정 체인을 둔 형태이다.
  \item 이 구조는 순수 중앙집중형 \texttt{CGW}가 아니라, \texttt{domain-controller style + CGW supervision}에 가깝다.
  \item 도메인 ECU는 각자 상태와 기능 의미를 유지하고, \texttt{V2X}와 \texttt{ADAS}가 경고 의미를 정규화·중재하며, \texttt{CGW}는 경계·route·fail-safe를 감독한다.
\end{itemize}
\end{frame}

\begin{frame}{핵심 용어와 ECU 약어 정의}
\begin{itemize}
  \item \texttt{CGW}: 차량 내부 경계와 route/fail-safe를 관리하는 gateway
  \item \texttt{V2X}: 외부 긴급 이벤트를 정규화하는 ingress owner
  \item \texttt{ADAS}: TTC와 긴급 접근을 결합해 최종 경고를 선택하는 decision owner
  \item \texttt{BCM / CLU / HUD / AMP}: 선택된 경고를 조명, 계기, 표시, 오디오로 반영하는 output surface
  \item \texttt{TEST\_SCN / TEST\_BAS}: 시나리오 주입과 verdict 집계를 담당하는 validation harness
\end{itemize}
\end{frame}

\begin{frame}{ECU 아키텍처 분류 기준}
\begin{itemize}
  \item 본 프로젝트는 ECU를 기능 이름이 아니라 역할 기준으로 분류하였다.
  \item 기준은 네 가지였다: Gateway/Backbone, Domain Runtime Owner, Local Feature/Output Surface, Validation Harness.
  \item 이 분류를 통해 누가 기능 의미를 소유하고, 누가 경계와 검증만 담당하는지를 명확히 했다.
\end{itemize}
\end{frame}

\begin{frame}{ECU 역할 분류}
\begin{itemize}
  \item \textbf{Gateway / Backbone}: 도메인 간 상태를 넘기고 경계 조건과 route를 감독하는 계층
  \item \textbf{Domain Runtime Owner}: 각 기능 의미를 직접 계산하고 최종 상태를 소유하는 계층
  \item \textbf{Local Feature / Output Surface}: 선택된 결과를 차체 기능이나 운전자 채널에 반영하는 계층
  \item \textbf{Validation Harness}: 시나리오 입력과 관측을 분리해 시험 기준을 유지하는 계층
\end{itemize}
\end{frame}

\begin{frame}{핵심 ECU 구성}
\begin{itemize}
  \item 입력 수집과 전달 경계는 \texttt{CGW}, \texttt{ETH\_BACKBONE}, \texttt{IVI}, \texttt{SGW}, \texttt{DCM}이 담당하였다.
  \item 핵심 판단은 \texttt{V2X}와 \texttt{ADAS}가 담당하고, 최종 출력은 \texttt{BCM}, \texttt{CLU}, \texttt{HUD}, \texttt{AMP}가 수행한다.
  \item \texttt{VALIDATION\_HARNESS}는 테스트 시나리오 주입과 결과 관측을 위해 별도 계층으로 유지하였다.
\end{itemize}
\end{frame}

\begin{frame}{핵심 경고 체인에서의 ECU 역할}
\begin{itemize}
  \item \texttt{CGW}는 차량 내부 상태를 모으고 판단 계층으로 넘기는 입력 관문 역할을 맡는다.
  \item \texttt{V2X}는 외부 긴급 이벤트를 정규화해 비교 가능한 입력으로 만든다.
  \item \texttt{ADAS}는 긴급 접근 정보와 객체 위험을 결합해 최종 경고와 감속 보조 요청을 결정한다.
  \item \texttt{BCM}, \texttt{CLU}, \texttt{HUD}, \texttt{AMP}는 선택된 경고를 조명, 계기, 전면 표시, 오디오 채널에 반영한다.
\end{itemize}
\end{frame}

\begin{frame}{도메인별 ECU 배치}
\begin{itemize}
  \item \textbf{Integration}: \texttt{CGW}, \texttt{ETHB}, \texttt{DCM}, \texttt{IBOX}, \texttt{SGW}
  \item \textbf{Powertrain}: \texttt{EMS}, \texttt{TCU}, \texttt{VCU}, \texttt{BAT\_BMS}, \texttt{OBC}, \texttt{DCDC}, \texttt{MCU}, \texttt{INVERTER}
  \item \textbf{Chassis / Safety}: \texttt{ESC}, \texttt{MDPS}, \texttt{ABS}, \texttt{EPB}, \texttt{EHB}, \texttt{VSM}, \texttt{RWS}
  \item \textbf{Body / Comfort}: \texttt{BCM}, \texttt{DATC}, \texttt{AFLS}, \texttt{AHLS}, \texttt{DOOR\_*}, \texttt{SEAT\_*}
  \item \textbf{IVI / HMI}: \texttt{IVI}, \texttt{CLU}, \texttt{HUD}, \texttt{TMU}, \texttt{AMP}, \texttt{NAV}, \texttt{DKEY}
  \item \textbf{ADAS / V2X}: \texttt{ADAS}, \texttt{V2X}, \texttt{SCC}, \texttt{LDWS\_LKAS}, \texttt{FCA}, \texttt{BCW}, \texttt{LCA}, \texttt{AVM}, \texttt{DMS}, \texttt{OMS}
\end{itemize}
\end{frame}

\subsection{네트워크와 통신 설계}
\begin{frame}{입력 계층에서 수집한 정보}
\begin{itemize}
  \item 차량 속도, 주행 상태, 조향 입력, 페달 상태와 같은 차량 기본 상태
  \item 구간 타입, 방향, 거리, 제한속도와 같은 구간 맥락 정보
  \item 긴급차량 유형, 방향, ETA, SourceID와 같은 V2X 이벤트 정보
  \item 객체 거리, 상대속도, 신뢰도와 같은 위험 판단 정보
\end{itemize}
\end{frame}

\begin{frame}{출력 계층에서 반영한 정보}
\begin{itemize}
  \item 앰비언트 채널은 색상, 패턴, 강조 모드로 경고 수준을 표현한다.
  \item 클러스터와 전면 표시 채널은 경고 문구, 방향, 팝업, 테마 상태를 반영한다.
  \item 오디오 채널은 포커스와 감쇄 정책을 조정하고, 제어 채널은 감속 보조 요청을 전달한다.
\end{itemize}
\end{frame}

\begin{frame}{도메인 CAN + Backbone 채택 배경}
\begin{itemize}
  \item 차량 상태, 구간 맥락, 차체·동력 기본 루프는 각 도메인 CAN에서 유지하는 것이 가장 자연스럽다.
  \item 반면 긴급 이벤트, 최종 경고 결과, fail-safe 상태는 cross-domain 의미이므로 backbone seam으로 올려야 설명과 검증이 쉬워진다.
  \item 본 설계의 핵심은 Ethernet 자체보다, \textbf{domain-local state collection + backbone-level warning meaning exchange}를 분리한 데 있다.
\end{itemize}
\end{frame}

\begin{frame}{Seam 관점의 설계 기준}
\begin{itemize}
  \item 본 프로젝트는 메시지보다 먼저 seam을 정의하고, seam마다 \textit{Owner / Bus / Timeout authority / Route authority}를 지정하였다.
  \item 예를 들어 Navigation context는 \texttt{IVI}가 owner이고, Emergency context는 \texttt{V2X}, Arbitration decision은 \texttt{ADAS}, Effective result는 \texttt{CGW}가 담당한다.
  \item 이 기준 덕분에 경고 의미, 경계 정책, 복귀 정책을 transport와 분리해서 설명할 수 있었다.
\end{itemize}
\end{frame}

\begin{frame}{CAN 기반 경로}
\begin{itemize}
  \item CAN은 차량 상태, 구간 정보, 차체 상태, 인포테인먼트 상태와 같은 기본 입력을 수집하는 주 경로로 사용되었다.
  \item 게이트웨이는 이 입력을 정규화하여 경고 판단 계층이 공통 형식으로 사용할 수 있도록 전달하였다.
  \item 따라서 CAN 경로의 역할은 세부 기능 수행보다 상태 수집과 기준선 형성에 가깝다.
\end{itemize}
\end{frame}

\begin{frame}{Backbone 기반 경로}
\begin{itemize}
  \item backbone 경로는 긴급 이벤트, 최종 경고 결과, fail-safe 상태처럼 cross-domain 의미를 전달하는 데 사용되었다.
  \item 핵심은 물리 매체보다도, raw signal이 아니라 정규화된 의미 seam만 backbone에 올린다는 점이었다.
  \item 이 경로를 통해 긴급 이벤트 선택과 최종 경고 상태를 시스템 수준에서 일관되게 관찰할 수 있었다.
\end{itemize}
\end{frame}

\begin{frame}{중요 통신 설계 원칙}
\begin{itemize}
  \item 기능 의미와 전송 수단을 먼저 분리하고, 이후 CAN과 Ethernet 경계를 연결하였다.
  \item 경고 의미의 owner는 기능 ECU가 유지하고, 검증 harness는 외생 입력과 관측 역할만 담당하도록 구분하였다.
  \item 운영 경로와 관측 경로를 함께 설계하여, 기능 동작과 시험 증거가 서로 충돌하지 않도록 했다.
\end{itemize}
\end{frame}

\begin{frame}{CAN ID 배정 구조}
\begin{itemize}
  \item 현재 active execution은 \texttt{11-bit CAN ID} 기준으로 운영한다.
  \item \texttt{0x100 \textasciitilde 0x13F}: 차량 기본 상태, 조향, 제동, 동력 기본 루프
  \item \texttt{0x1C0 \textasciitilde 0x20F}: 긴급 알림, 위험 판단, 주행 보조 상태
  \item \texttt{0x260 \textasciitilde 0x27F}: 출입, 조명, 공조, 실내 편의 상태 / \texttt{0x280 \textasciitilde 0x29F}: 내비 문맥, 클러스터, HMI 상태
  \item backbone 식별자는 별도로 관리하며, logical contract 값으로 \texttt{0x510 \textasciitilde 0x512}, \texttt{0xE100}, \texttt{0xE200}을 사용한다.
\end{itemize}
\end{frame}

\begin{frame}{대표 메시지 계약}
\begin{itemize}
  \item \texttt{frmVehicleStateCanMsg (0x2A0)}: \texttt{VCU}가 차량 기본 상태를 제공하는 입력 계약
  \item \texttt{frmNavigationRouteMsg (0x285)}: \texttt{NAV}가 구간 문맥과 경로 정보를 제공하는 입력 계약
  \item \texttt{ETH\_EmergencyAlert (0xE100)}: \texttt{V2X}가 긴급 이벤트를 수신하는 backbone ingress 계약
  \item \texttt{ethSelectedAlertMsg (0xE200)}: \texttt{ADAS}가 선택된 경고 결과를 전달하는 backbone 출력 계약
  \item \texttt{ethFailSafeStateMsg (0xE212)}: \texttt{CGW}가 경계 건강도와 fail-safe 상태를 전달하는 계약
\end{itemize}
\end{frame}

\begin{frame}{핵심 네트워크 흐름}
\begin{itemize}
  \item 차량 기본 상태는 \texttt{VCU/MDPS/ESC}에서 CAN으로 생성되고, \texttt{CGW} 또는 domain owner가 공통 판단에 필요한 상태로 정규화한다.
  \item 구간 문맥은 \texttt{NAV/IVI}에서 CAN으로 유지되고, cross-domain이 필요한 경우 backbone seam으로 전달된다.
  \item 긴급 이벤트는 \texttt{V2X}가 \texttt{ETH\_EmergencyAlert}로 수신한 뒤, \texttt{Core/CoreState} ingress seam으로 정규화한다.
  \item 최종 경고와 fail-safe 결과는 \texttt{ADAS -> CGW -> BCM/CLU/HUD/AMP} 흐름으로 이어지며, backbone은 의미 전달 경로로만 사용된다.
\end{itemize}
\end{frame}

\begin{frame}{멀티버스와 가시성 정책}
\begin{itemize}
  \item 모든 ECU를 다중 버스에 붙이는 것이 아니라, foreign-domain CAN 가시성이 필요한 ECU만 추가 visibility를 부여하였다.
  \item \texttt{CGW}, \texttt{TEST\_SCN}은 true multibus anchor이고, \texttt{ADAS}, \texttt{BCM}, \texttt{IVI}, \texttt{VCU}는 필요한 foreign-domain CAN만 추가로 본다.
  \item \texttt{TEST\_BAS}는 의도적으로 single-bus observer로 유지하여, system-wide 의미와 transport 복잡도를 분리하였다.
\end{itemize}
\end{frame}

\begin{frame}{아키텍처 채택 근거}
\begin{itemize}
  \item 핵심 요구는 개별 경고 기능을 늘리는 것이 아니라, 서로 다른 주행 맥락을 하나의 선택 규칙과 복귀 규칙으로 묶는 것이었다.
  \item 이를 위해 \texttt{V2X}는 외부 이벤트를 정규화하고, \texttt{ADAS}는 위험을 중재하며, \texttt{CGW}는 경계·route·fail-safe를 담당하는 구조가 필요했다.
  \item 즉 채택안의 의미는 transport 변경 자체가 아니라, \textbf{owner 분리}, \textbf{boundary supervision}, \textbf{cross-domain seam 설명 가능성}을 동시에 확보한 데 있다.
  \item 이 기준 덕분에 구간 경고, 긴급차량 접근, 교차로·합류 위험을 동일한 시스템 규칙 위에서 설명할 수 있었다.
\end{itemize}
\end{frame}

\subsection{구현 전환과 운영 구조}
\begin{frame}{구현 시 중요했던 원칙}
\begin{itemize}
  \item 설계 의미와 transport를 먼저 분리하고, 이후에 CAN과 Ethernet 경계를 붙이는 방식으로 진행하였다.
  \item 기능 owner와 validation harness를 분리하여, 테스트 경로가 기능 의미를 직접 대체하지 않도록 유지하였다.
  \item \texttt{src/capl}과 active mirror를 함께 관리하며 compile과 runtime을 별도 게이트로 취급하였다.
\end{itemize}
\end{frame}

\begin{frame}{주요 설계 전환점}
\begin{itemize}
  \item 초기의 ``중앙 게이트웨이가 모든 raw payload를 relay한다''는 해석에서, \texttt{domain-controller style + CGW supervision} 구조로 관점을 바꾸었다.
  \item 메시지 단위 예외 처리 대신 \texttt{owner-route-layer separation} 규칙을 세워, seam마다 owner·timeout authority·route authority를 고정하였다.
  \item CAN ID는 29-bit 정책과 \texttt{11-bit active execution}을 분리해 문서와 실행 기준이 서로 충돌하지 않도록 정리하였다.
  \item 검증 경로도 운영 경로와 분리하여, \texttt{TEST\_SCN/TEST\_BAS}, \texttt{Test Matrix}, \texttt{EXT\_DIAG}가 관측과 verdict 역할만 수행하도록 재구성하였다.
\end{itemize}
\end{frame}

\begin{frame}{운영 경로와 검증 관측 경로}
\begin{itemize}
  \item 운영 경로는 실제 경고 의미를 계산하고 운전자 채널로 반영하는 경로이다.
  \item 검증 관측 경로는 상태, 경계 조건, 판정 결과를 독립적으로 확인하는 경로이다.
  \item 두 경로를 함께 설계함으로써 기능 동작과 시험 증거가 서로 충돌하지 않도록 하였다.
\end{itemize}
\end{frame}

\begin{frame}{진단·관측 계층}
\begin{itemize}
  \item \texttt{SGW}, \texttt{DCM}, \texttt{EXT\_DIAG}는 서비스 상태와 경계 조건을 독립적으로 관측하는 진단·관측 계층으로 구성하였다.
  \item 이 계층은 핵심 경고 로직을 직접 결정하지 않고, 상태 해석과 Fail-Safe 판단에 필요한 정보를 제공한다.
  \item 결과적으로 경고 판단 계층과 진단 관측 계층이 분리되어 시스템 설명이 더 명확해졌다.
\end{itemize}
\end{frame}

\begin{frame}{패널과 출력 채널의 설계 구성과 역할}
\begin{itemize}
  \item \texttt{Ambient}, \texttt{Cluster}, \texttt{Navigation}, \texttt{Diagnostic Console} 패널은 각각 경고 출력, 상태 확인, 맥락 표시, 진단 관측을 담당하는 출력화면이다.
  \item 패널은 독립적으로 판단하는 UI가 아니라, 정해진 입력 경로와 상태 확인 경로를 통해 결과만 일관되게 보여주도록 구성하였다.
  \item 따라서 패널 구성은 기능 결과가 어떤 채널로 노출되고 어떤 상태가 검증 가능한지를 보여주는 설계 근거가 되었다.
\end{itemize}
\end{frame}

\begin{frame}{출력 패널의 구성과 시각화}
\begin{itemize}
  \item 출력 패널은 각 시나리오에서 \texttt{Ambient}, \texttt{Cluster}, \texttt{Navigation}의 상태 변화가 어떻게 드러나는지를 시각적으로 확인하는 테스트 수단으로 활용되었다.
  \item 동일 시나리오에서 경고 레벨, 구간 맥락, 표시 채널 변화가 함께 보이도록 배치하여 출력 변화와 관측 결과를 한 화면에서 비교할 수 있게 하였다.
  \item 패널 작업의 핵심은 화면 변경보다 CAPL과 시스템 변수 계약에 맞춰 출력 의미가 일관되게 보이도록 정렬하는 데 있었다.
\end{itemize}
\end{frame}

\begin{frame}{데모 1: 전체 차량 패널}
\begin{itemize}
  \item 이 데모에서는 기본주행상태, 구간 전환, 경고 출력 채널이 하나의 흐름으로 연결되는 장면을 보여준다.
  \item 특히 속도·기어·구간 문맥이 바뀔 때 앰비언트, 클러스터, 전면 표시가 어떻게 함께 반응하는지 확인한다.
  \item 이 장면은 ``차량 전체 기준선''을 먼저 보여주고, 이후 V2X와 ADAS가 이 기준선 위에서 어떤 판단을 추가하는지 이해하게 해준다.
\end{itemize}
\end{frame}

\section{핵심 로직과 구현}
\subsection{V2X 기반 긴급 이벤트 정규화}

\begin{frame}{V2X 기능 역할}
\begin{itemize}
  \item V2X의 역할은 긴급 이벤트를 수신하는 데서 끝나지 않는다.
  \item 다수의 긴급 후보가 존재하더라도 ADAS와 HMI가 참조할 활성 이벤트 하나를 정규화하는 것이 핵심 역할이다.
  \item 즉 V2X는 외부 이벤트를 비교 가능한 입력으로 바꾸는 전처리 계층이다.
\end{itemize}
\end{frame}

\begin{frame}{V2X 선택 규칙}
\begin{block}{긴급 이벤트 선택 규칙}
\[
e_t^*=\arg\min(-rank(type), ETA, SourceID)
\]
\end{block}
\begin{itemize}
  \item 구급차가 경찰차보다 우선한다.
  \item 같은 유형이면 ETA가 더 작은 이벤트를 우선한다.
  \item ETA가 같으면 SourceID가 작은 이벤트를 선택한다.
\end{itemize}
\end{frame}

\begin{frame}{V2X timeout clear}
\begin{itemize}
  \item 긴급 이벤트는 한 번 들어왔다고 계속 유지되면 안 된다.
  \item 따라서 1000ms 이상 갱신되지 않으면 \texttt{timeout clear}를 발생시켜 안전하게 경고를 해제하도록 구성하였다.
  \item 이 규칙은 긴급 이벤트 유지와 종료를 시스템 수준에서 일관되게 설명하는 핵심 기준이 되었다.
\end{itemize}
\end{frame}

\begin{frame}{데모 2: V2X 동작}
\begin{itemize}
  \item 이 데모에서는 경찰과 구급 이벤트가 들어올 때 \texttt{type > ETA > SourceID} 규칙으로 활성 이벤트가 어떻게 정해지는지 보여준다.
  \item 또한 ETA 변화와 timeout clear가 패널과 상태값에 어떤 식으로 반영되는지도 함께 확인한다.
  \item 즉 이 장면은 V2X가 ``비교 가능한 입력 정규화''를 수행하고, 그 결과가 다음 ADAS 판단의 직접 입력으로 넘어간다는 점을 보여준다.
\end{itemize}
\end{frame}

\subsection{ADAS 기반 위험도 산정과 경고 중재}
\begin{frame}{ADAS 기능 역할}
\begin{itemize}
  \item ADAS는 긴급 접근 정보와 객체 위험 정보를 결합해 최종 경고와 감속 보조 요청을 산출한다.
  \item 즉 ADAS는 단순 위험 계산기가 아니라, 다중 입력을 단일 경고 정책으로 수렴시키는 판단 계층이다.
  \item 이 계층에서 경고 선택, 경고 유형, 경고 레벨, 감속 보조 요청이 함께 결정된다.
\end{itemize}
\end{frame}

\begin{frame}{긴급 접근 위험도}
\begin{block}{Emergency Risk}
\[
R_{emg}=0.55\cdot etaScore + speedScore + directionScore + typeBias
\]
\end{block}
\begin{itemize}
  \item ETA가 짧을수록, 자차 속도가 높을수록, 접근 방향이 직접적일수록 위험도가 높아진다.
  \item 긴급차량 유형에 따른 bias를 더해 같은 접근이라도 경고 우선순위 차이를 반영하였다.
\end{itemize}
\end{frame}

\begin{frame}{객체 TTC 계산}
\small
\begin{block}{Time-to-Collision}
\[
TTC=\frac{3600\cdot d}{|v_{rel}|}
\]
\end{block}
\begin{itemize}
  \item \(d\)는 상대 거리, \(v_{rel}\)은 상대 속도이다.
  \item TTC는 교차로 측방 접근과 합류구간 위험을 하나의 시간 축에서 비교하는 핵심 기준이다.
\end{itemize}
\end{frame}

\begin{frame}{객체 위험도}
\small
\begin{itemize}
  \item ADAS는 TTC만 보지 않고 상대 거리, 객체 신뢰도, 시나리오 bias를 함께 반영해 객체 위험도를 계산한다.
  \item 그 결과 객체 위험은 교차로, 합류, 급간섭 상황을 하나의 위험도 축으로 정리할 수 있게 되었다.
  \item 이는 긴급 접근 정보와 TTC 위험을 동일한 중재 구조 안에서 다룰 수 있게 만든 핵심 단계였다.
\end{itemize}
\end{frame}

\begin{frame}{최종 위험도 선택}
\small
\begin{itemize}
  \item 객체 위험도는 \texttt{baseRisk + rangeBias + confidenceBias + scenarioBias}로 구성되고, 긴급 접근 위험도와 함께 최종 위험도 선택 단계로 전달된다.
  \item 최종 위험도는 \(R_t = \max(R_{emg}, R_{obj})\)로 정의하여, 긴급 접근 위험과 객체 위험 중 더 높은 쪽이 우선 반영되도록 하였다.
  \item 이 단계에서 긴급 접근과 객체 위험은 하나의 경고 정책 안에서 비교 가능한 동일 판단 축으로 통합된다.
\end{itemize}
\end{frame}

\begin{frame}{계층형 경고 중재}
\small
\begin{itemize}
  \item V2X가 정규화한 긴급 이벤트와 ADAS가 계산한 객체 위험은 최종적으로 하나의 경고 상태로 중재된다.
  \item 이 과정에서 긴급 경고, 구간 경고, 객체 위험 경고가 동시에 유효하더라도 최종 경고는 단일 상태로 결정된다.
  \item 따라서 운전자는 현재 상황에서 가장 중요한 경고를 일관되게 인지할 수 있다.
\end{itemize}
\end{frame}

\begin{frame}{감속 보조 on/off 안정화}
\small
\begin{itemize}
  \item 감속 보조는 위험도가 임계값을 넘는다고 즉시 반복 진동해서는 안 된다.
  \item 따라서 on 조건과 off 조건을 분리한 hysteresis로 감속 보조 상태를 안정화하였다.
  \item 이 구조는 교차로와 합류구간에서 불필요한 감속 요청 반복을 줄이는 데 중요했다.
\end{itemize}
\end{frame}

\begin{frame}{fail-safe와 복귀 규칙}
\begin{itemize}
  \item 전달 경계 이상이나 stale 입력이 발생하면 시스템은 fail-safe 상태로 강등되어야 한다.
  \item 이때 \texttt{CGW}는 \texttt{warningPathStatus}, \texttt{failSafeMode}, \texttt{selectedAlertEffective*}를 함께 shaping한다.
  \item 정상 경로가 복구되면 최소 경고를 해제하고 정상 effective alert 경로로 안전하게 복귀해야 한다.
\end{itemize}
\end{frame}

\begin{frame}{데모 3: ADAS 위험도와 경고}
\begin{itemize}
  \item 이 데모에서는 객체 TTC 변화, 교차로·합류구간 위험, 감속 보조 요청이 실제로 어떤 순서로 반영되는지 보여준다.
  \item 긴급 접근 정보와 객체 위험이 동시에 존재할 때 최종 경고가 하나로 수렴하는 장면도 함께 확인한다.
  \item 이 장면은 ADAS가 위험도와 복귀를 함께 다루는 중재 계층임을 보여주고, 이어지는 시험 체계가 정확히 이 로직을 어떻게 검증했는지로 자연스럽게 연결된다.
\end{itemize}
\end{frame}

\subsection{CAPL 구현 구조}
\begin{frame}{CAPL 기반 구현 구조}
\begin{itemize}
  \item 구현은 CAPL 기반 논리 ECU를 중심으로 수행하였다.
  \item \texttt{V2X}, \texttt{ADAS}, \texttt{CGW}, \texttt{IVI}, \texttt{BCM}, \texttt{CLU}, \texttt{HUD} 등은 각자 입력, 판단, 출력 역할을 분담하였다.
  \item 이를 통해 기능 구조와 코드 구조가 직접 대응되도록 유지하였다.
\end{itemize}
\end{frame}

\begin{frame}{구현 단계에서 확인한 구조적 원칙}
\begin{itemize}
  \item 기능 로직이 깊어질수록 설계 문서와 시험 기준을 같이 유지하지 않으면 전체 구조가 쉽게 흔들린다.
  \item CAN과 Ethernet이 혼재된 구조에서는 ownership과 observation을 명시하지 않으면 디버깅이 어려워진다.
  \item 따라서 구현은 코드 작성보다 경계와 의미를 정리하는 작업에 더 가까웠다.
\end{itemize}
\end{frame}

\section{검증 체계와 운영 자동화}
\subsection{검증 설계와 관측 구조}

\begin{frame}{Validation Harness의 역할}
\begin{itemize}
  \item \texttt{TEST\_SCN}은 시나리오 입력을 담당하고, \texttt{TEST\_BAS}는 결과 관측과 baseline 판정을 담당한다.
  \item 핵심 경고 의미는 기능 ECU가 소유하고, harness는 이를 대신 생성하지 않는다는 원칙을 유지하였다.
  \item 이 분리는 시험 자산이 기능 로직을 오염시키지 않도록 하는 핵심 기준이었다.
\end{itemize}
\end{frame}

\begin{frame}{Diagnostic Console의 의미}
\begin{itemize}
  \item Diagnostic Console은 서비스 상태, 경계 조건, 진단 상태를 독립적으로 확인하는 시험 관측면으로 활용되었다.
  \item 이를 통해 경고 판단 문제와 서비스/경계 문제를 분리해 해석할 수 있었다.
  \item 즉 Diagnostic Console은 기능 화면이 아니라, 시험 설계와 결과 해석을 돕는 관측 도구에 가깝다.
\end{itemize}
\end{frame}

\begin{frame}{시험 전략}
\begin{itemize}
  \item 시험은 앞서 설명한 V2X 정규화, ADAS 중재, 복귀 규칙을 검증하기 위해 단위시험, 통합시험, 시스템시험의 세 층으로 구성하였다.
  \item 단위시험은 로직 규칙을, 통합시험은 중재와 복귀를, 시스템시험은 연속 시나리오와 fail-safe를 다루도록 정리하였다.
  \item 이렇게 함으로써 같은 경고 컨셉을 서로 다른 수준에서 반복 확인할 수 있었다.
\end{itemize}
\end{frame}

\subsection{시험 자산과 실행 체계}
\begin{frame}{단위시험 검증 항목}
\begin{itemize}
  \item 긴급 이벤트 수신, 유지, timeout clear
  \item TTC 기반 위험도 산정과 객체 위험 분류
  \item 감속 보조 요청과 운전자 개입 해제
  \item 채널 가용성, fail-safe, 최소 경고 유지
\end{itemize}
\end{frame}

\begin{frame}{통합시험 검증 항목}
\begin{itemize}
  \item 경찰과 구급의 동시 접근 시 우선순위 규칙
  \item ETA와 SourceID 동률 처리
  \item 긴급 경고와 TTC 위험이 동시에 유효한 경우의 최종 경고 선택
  \item timeout clear 이후 해제와 복귀의 일관성
\end{itemize}
\end{frame}

\begin{frame}{시스템시험 검증 항목}
\begin{itemize}
  \item 외부 송신 주기 유지
  \item 연속 시나리오에서 경고 선택과 복귀
  \item fail-safe 진입과 정상 복귀
  \item 운영 경로와 관측 경로가 함께 유지되는지 여부
\end{itemize}
\end{frame}

\begin{frame}{대표 시험 시나리오}
\begin{itemize}
  \item 구간 경고와 긴급 접근이 동시에 유효한 상황에서 최종 경고가 하나로 수렴하는지 확인하였다.
  \item 경찰과 구급의 동시 접근, ETA 및 SourceID 동률 해소, timeout clear 이후 복귀를 핵심 규칙 검증 항목으로 유지하였다.
  \item 교차로·합류구간 TTC 위험과 fail-safe 복귀는 시스템시험 시나리오에서 별도로 반복 확인하였다.
\end{itemize}
\end{frame}

\begin{frame}{대표 시나리오 end-to-end}
\begin{itemize}
  \item 대표 사례는 ``긴급차량 접근 + 교차로 또는 합류구간 위험''이 동시에 유효한 상황이다.
  \item 입력은 V2X 이벤트와 객체 위험으로 들어오고, V2X는 활성 이벤트를 정규화하며 ADAS는 이를 TTC 기반 위험과 결합해 최종 경고와 감속 보조 요청을 결정한다.
  \item 이후 결과는 앰비언트, 클러스터, 전면 표시, 오디오, 감속 보조 채널로 동시에 반영되고, 같은 시나리오가 단위시험·통합시험·시스템시험으로 이어진다.
\end{itemize}
\end{frame}

\begin{frame}{실행 증거와 관측 화면}
\begin{itemize}
  \item 실행 증거는 단순 PASS 수가 아니라, \texttt{01 Requirements}, \texttt{03 Function}, \texttt{05/06/07} 시험 자산이 같은 ID 기준으로 연결되어 있음을 보여주는 자료다.
  \item CANoe Test Unit과 Trace 창에서는 각 시험 ID별 PASS/FAIL과 판정 근거를 확인할 수 있고, 결과는 엑셀 및 보고서 형태로 다시 정리된다.
  \item 따라서 실행 화면과 보고서는 ``기능이 동작했다''는 증거이면서 동시에 ``요구사항부터 시험 결과까지 추적된다''는 증거이기도 하다.
\end{itemize}
\end{frame}

\begin{frame}{데모: 시험 실행과 리포트}
\begin{itemize}
  \item 이 데모에서는 Trace 창, Test Unit 실행 창, 상세 리포트를 순서대로 보여주며 시험 구조가 실제로 동작함을 설명한다.
  \item Trace는 seam과 메시지 흐름을, Test Unit 창은 시험 ID별 PASS/FAIL 상태를, 리포트는 시나리오 단위 verdict와 근거를 보여준다.
  \item 이 장면은 요구사항, 기능 정의, 시험 자산, 결과 보고서가 같은 기준선 위에서 운영된다는 점을 가장 직접적으로 보여주는 증거이다.
\end{itemize}
\end{frame}

\begin{frame}{추적과 리포트 구조}
\begin{itemize}
  \item 요구사항, HARA, 기능 정의를 기준으로 검증 계약, Test Oracle, Test Matrix, 단위시험·통합시험·시스템시험 자산을 구성하였다.
  \item 그 결과 시험 항목 번호만 남는 것이 아니라, 어떤 기능을 어떤 기준으로 검증했고 어떤 보고서로 귀결되는지가 함께 남도록 정리할 수 있었다.
  \item 이 구조는 00--07 문서, 시험 실행 결과, 엑셀 산출을 하나의 추적선 위에서 유지하는 데 핵심 역할을 했다.
\end{itemize}
\end{frame}

\begin{frame}{운영 자동화 제품의 역할}
\begin{itemize}
  \item 운영 자동화 제품은 시험 실행 이후 결과 수집, XML 기반 집계, 문서 closeout, 엑셀 산출을 같은 흐름에서 처리하도록 구성되었다.
  \item 이를 통해 개별 시험 결과를 사람이 수작업으로 옮기는 대신, 기준선 유지와 결과 정리에 필요한 반복 작업을 줄일 수 있었다.
  \item 따라서 운영 자동화는 부가 기능이 아니라, 정량 결과를 안정적으로 확보하기 위한 실무 기반으로 작동한다.
\end{itemize}
\end{frame}

\section{결과와 해석}
\subsection{핵심 결과와 정량 지표}

\begin{frame}{프로젝트 성과 요약}
\begin{itemize}
  \item 본 프로젝트는 개별 기능 시연이 아니라, 전체 차량 기준선 위에서 구간 경고, 긴급차량 접근, 객체 위험을 하나의 경고 정책으로 통합한 시스템 구축 결과다.
  \item 이 구조는 \texttt{V2X}, \texttt{ADAS}, \texttt{CGW} 논리 ECU의 CAPL 구현과 출력·진단·시험 경로까지 포함하여 실제 실행 가능한 기준선으로 정리되었다.
  \item 또한 요구사항, 기능 정의, 흐름, 통신, 변수, 코드, 단위시험·통합시험·시스템시험이 하나의 추적 구조로 닫히도록 구성하였다.
\end{itemize}
\end{frame}

\begin{frame}{핵심 정량 지표 정의}
\begin{itemize}
  \item 결과는 세 가지 핵심 지표를 기준으로 집계하도록 정의하였다.
\end{itemize}
\begin{block}{정량 지표}
\[
\begin{aligned}
N1&=(S_{sel}/T_{v2x})\times100 && \text{V2X 긴급 이벤트 선택 정확도}\\
N2&=(S_{int}/T_{int})\times100 && \text{교차로·합류구간 통합 경고 성공률}\\
N3&=(S_{rec}/T_{rec})\times100 && \text{복귀 안정성 준수율}
\end{aligned}
\]
\end{block}
\end{frame}

\begin{frame}{정량 지표와 데이터 출처}
\begin{itemize}
  \item \textbf{N1}: emergencyType, ETA, SourceID, selected alert verdict를 기준으로 집계한다.
  \item \textbf{N2}: objectTtcMin, objectRiskClass, decelAssistReq, selected alert verdict를 기준으로 집계한다.
  \item \textbf{N3}: timeoutClear, warningPathStatus, effective alert state, baseline verdict를 기준으로 집계한다.
  \item 핵심은 단순 PASS 수가 아니라, owner-published seam과 시나리오 verdict를 함께 읽어 기능 의미를 평가하는 데 있다.
\end{itemize}
\end{frame}

\begin{frame}{왜 이 세 지표가 중요한가}
\begin{itemize}
  \item \(N1\)은 외부 이벤트를 비교 가능한 입력으로 정규화하는 능력을 보여준다.
  \item \(N2\)는 긴급 이벤트와 TTC 위험을 하나의 경고 정책으로 수렴시키는 능력을 보여준다.
  \item \(N3\)는 시스템이 경고를 어떻게 안전하게 해제하고 복귀시키는지를 보여준다.
\end{itemize}
\end{frame}

\begin{frame}{핵심 결과 해석}
\begin{itemize}
  \item 본 프로젝트의 강점은 단일 시나리오를 설명하는 수준을 넘어, 전체 차량 시스템을 처음부터 끝까지 하나의 기준선으로 연결했다는 점이다.
  \item 특히 77개 단위시험, 45개 통합시험, 46개 시스템시험 자산은 핵심 기능이 요구사항부터 시스템시험까지 이어지는 검증 구조를 갖추었음을 보여준다.
  \item 따라서 결과는 단순 PASS 개수보다 V2X 선택, 통합 경고, 복귀 안정성을 어떤 근거와 시나리오로 해석할 수 있는지가 더 중요하다.
\end{itemize}
\end{frame}

\section{프로젝트 정리 및 향후 계획}
\subsection{프로젝트 가치와 향후 방향}

\begin{frame}{프로젝트 성과와 구축 범위}
\begin{itemize}
  \item 구간 경고, 긴급차량 접근, 교차로·합류구간 위험을 분리된 기능으로 남기지 않고, 전체 차량 기준선 위에서 하나의 통합 경고 정책으로 구현하였다.
  \item 문서, CAPL ECU 코드, 시험 자산, 운영 자동화 산출을 같은 기준선 위에 정렬하여 확장과 검증이 가능한 구조를 확보하였다.
  \item 이 점에서 본 프로젝트는 단순 기능 시연이 아니라 차량 시스템 관점의 통합 구축 사례로 의미를 가진다.
\end{itemize}
\end{frame}

\begin{frame}{프로젝트 확장성과 활용 관점}
\begin{itemize}
  \item 실제 차량 시스템에서는 입력이 많아질수록 경고 기능도 늘어나지만, 사용자가 경험하는 것은 결국 하나의 최종 경고 정책이다.
  \item 본 프로젝트는 이 정책을 문서, ECU 구조, 시험 자산, 운영 체계까지 포함한 형태로 설명할 수 있게 했다는 점에서 의미가 있다.
  \item 특히 향후 SDV 환경에서는 기능 추가보다 정책 통합, 관측 가능성, 반복 가능한 검증 구조가 더 중요해지므로, 본 프로젝트의 경험은 그 출발점으로 활용될 수 있다.
\end{itemize}
\end{frame}

\begin{frame}{프로젝트를 통해 배운 점}
\begin{itemize}
  \item 이번 프로젝트를 하며 가장 크게 느낀 점은, 서로 다른 입력을 하나의 경고 정책으로 정리하는 구조를 세우는 일이었다.
  \item 설계, 구현, 시험, 운영 자동화는 따로 최적화할 수 없고 같은 추적 구조 안에서 함께 움직여야 안정적으로 유지된다는 점을 배웠다.
  \item 혼합 CAN/Ethernet 경로에서는 기능 경로와 관측 경로를 함께 설계해야 실제 디버깅과 검증이 가능했다.
  \item 교육생 입장에서 특히 크게 배운 부분은, 코드를 많이 쓰는 것보다 ECU 간 역할 경계와 메시지 의미를 끝까지 맞추는 일이 시스템 프로젝트에서 더 중요하다는 점이었다.
\end{itemize}
\end{frame}

\begin{frame}{향후 연구 계획}
\begin{itemize}
  \item 이미 구축된 XML 기반 결과 집계와 운영 자동화 기반을 활용해 반복 시험과 환경 통제 실험을 수행한다.
  \item 속도, ETA, TTC, timeout과 같은 주요 파라미터 변화와 ECU 기능 업데이트 전후 결과를 동일 기준선에서 비교한다.
  \item 이를 통해 정량 결과를 확정하고, 통계적으로 의미 있는 수치값과 변화 경향을 도출하여 결과 해석의 신뢰도를 높인다.
  \item 협력인지, SPaT/MAP, 도로위험·군집 이벤트 공유와 같은 확장형 V2X 주제 전개
\end{itemize}
\end{frame}

\begin{frame}{맺음말}
\begin{itemize}
  \item 본 프로젝트는 전체 차량 기준선 위에서 주행 경고 정책을 설계, 구현, 시험, 운영 기준선으로 끝까지 연결한 사례이다.
  \item 발표의 핵심은 단일 기능 소개보다, 요구사항부터 시스템시험까지 닫히는 시스템 통합 구조를 실제 ECU 구현과 함께 구축했다는 점에 있다.
  \item 이후 단계에서는 현재 확보한 기준선을 바탕으로 정량 결과를 확정하고 확장 기능을 점진적으로 누적할 수 있다.
\end{itemize}
\end{frame}

\end{document}
